\chapter{Python概述}
\section{Python 版本区别，Python 3 和 Python 2 区别详解}

\subsection{print 由语句变成函数}
Python 3.x \codeinline{print} 函数代替了 \codeinline{print} 语句

\begin{lstlisting}[style=demostyle,numbers=none]
# Python 2.x
print "Hello World"
print "3,4"

# Python 3.x
print("Hello World")
print("3,4")
\end{lstlisting}

\subsection{编码问题}
Python 3.x 默认使用 \codeinline{UTF-8} 编码，Python 2.x 默认使用 \codeinline{ASCII} 编码

\begin{lstlisting}[style=demostyle,numbers=none]
>>> import sys
>>> sys.getdefaultencoding()
'utf-8'
\end{lstlisting}

\subsection{整数除法}
Python 3.x 中整数除法返回浮点数，Python 2.x 中整数除法返回整数，需要使用 \codeinline{//}
\begin{lstlisting}[style=demostyle,numbers=none]
# Python 2.x
>>> 3/2
1
>>> 3.0/2
1.5

# Python 3.x
>>> 3/2
1.5
>>> 3.0/2
1.5
\end{lstlisting}        

\subsection{异常}
\begin{enumerate}[label=\arabic*)]
    \item 在 Python 2.x 版本中，所有类型的对象都是直接被抛出的，但是在 Python 3.x 版本中，只有继承 \codeinline{BaseException} 的对象才可以被抛出。
    \item 在 Python 2.x 版本中，捕获异常的语法是“\codeinline{except Exception, var:}”；但在 Python 3.x 版本中，引入了 \codeinline{as} 关键字，捕获异常的语法变更为 “\codeinline{except Exception as var:}”。
    \item 在 Python 3.x 版本中，处理异常用 “\codeinline{raise Exception(args)}”代替了“\codeinline{raise Exception, args}”。
    \item Python 3.x 版本中，取消了异常类的序列行为和 \codeinline{.message} 属性。
\end{enumerate}
\begin{lstlisting}[style=demostyle,numbers=none]
#Python 2.x
>>> try:
... raise TypeError,"类型错误" 
... except TypeError,err:
... print err.message
...
类型错误
#Python 3.x
>>> try:
... raise TypeError("类型错误")
... except TypeError as err:
... print(err)
...
类型错误
\end{lstlisting}   

\subsection{八进制字面量}
在 Python 3.x 中，表示八进制字面量的方式只有一种，并且必须写成“\codeinline{0o1000}”这样的方式，原来“\codeinline{01000}”的方式不能使用了。举个例子： 
\begin{lstlisting}[style=demostyle,numbers=none]
#Python 2.x
>>> 0o1000
512
>>> 01000
512
#Python 3.x
>>> 01000
 File "<stdin>", line 1
 01000
 ^ 
SyntaxError: invalid token
>>> 0o1000
512
\end{lstlisting}

\section{Python 2to3}
由于 Python 2.x 和 Python 3.x 的差别较大，因此 Python 2.x 的多数代码无法直接在 Python 3.x 环境中运行。
而由于兼容性的原因，我们在网络上查找的资源多数是 Python 2.x 的代码，如果想要在 Python 3.x 环境下运
行，就需要修改源代码。

\begin{lstlisting}[style=demostyle,numbers=none]
2to3 -w test.py # -w 选项表示直接修改源文件
\end{lstlisting}



